% Generated draft: requires mathematical and source review.
\documentclass[11pt]{article}
\usepackage[margin=25mm]{geometry}
\usepackage{fontspec}
\setmainfont{Latin Modern Roman}
\usepackage{amsmath,amssymb,mathtools,unicode-math}
\setmathfont{Latin Modern Math}
\usepackage{enumitem,needspace,xurl,hyperref}
\setlength{\parindent}{0pt}
\setlength{\parskip}{5pt}
\setlength{\emergencystretch}{3em}
\newcommand{\N}{\mathbb{N}}
\newcommand{\Z}{\mathbb{Z}}
\newcommand{\Q}{\mathbb{Q}}
\newcommand{\R}{\mathbb{R}}
\newcommand{\C}{\mathbb{C}}
\newcommand{\F}{\mathbb{F}}
\newcommand{\A}{\mathbb{A}}
\newcommand{\PP}{\mathbb{P}}
\newcommand{\E}{\mathbb{E}}
\newcommand{\eps}{\varepsilon}
\newcommand{\dd}{\,\mathrm{d}}
\begin{document}
\section*{UnsolvedMath ID: 30003545}
\textbf{Source problem code:} catalogue entry
\subsection{Definitions and mathematical statement}
Let \(\Sigma=S^{n-2}\subset S^n\) be the singular locus and let \(M=S^n\setminus\Sigma\). The source denotes by \(h_\beta\) the standard edge-cone metric on \((S^n,\Sigma)\) with transverse cone angle \(2\pi\beta\). A transverse cone of this angle has metric
\[
\dd r^2+\beta^2r^2\dd\theta^2,
\qquad \theta\in\R/(2\pi\Z).
\]
This specifies the cone-angle convention; the supplied record does not give the full metric \(h_\beta\), its precise asymptotic regularity, or the admissible conformal factors at \(\Sigma\).

Under the variational interpretation suggested by the supplied literature assessment, an edge-cone Yamabe metric means an admissible metric attaining the Yamabe infimum in the prescribed conformal class. More explicitly, for a fixed dimension \(n\geq3\), let \(\mathcal G_\beta\) denote the intended class of admissible edge-cone metrics
\[
g=u^{4/(n-2)}h_\beta,\qquad u>0\text{ on }M,
\]
with the prescribed cone angle and with finite positive volume and an integrable scalar curvature. Write
\[
\mathcal Y(g)=
\frac{\displaystyle\int_M\operatorname{Scal}_g\,\dd V_g}
{\displaystyle\operatorname{Vol}_g(M)^{(n-2)/n}},
\qquad
Y_\beta=\inf_{g\in\mathcal G_\beta}\mathcal Y(g).
\]
The proposed nonexistence statement is
\[
\forall\beta\in(1,2),\qquad
\nexists\,g\in\mathcal G_\beta
\quad\text{such that}\quad
\mathcal Y(g)=Y_\beta.
\]
Thus the conjectured nonexistence concerns cone angles strictly between \(2\pi\) and \(4\pi\); neither endpoint belongs to the requested range.

\emph{Unresolved qualifications.} This is a conditional variational formalization: the source does not specify the dimension range, define the full standard metric, or determine \(\mathcal G_\beta\), including regularity and boundary behavior at \(\Sigma\). The convention \(n\geq3\) is required for the displayed usual Yamabe formulation and is not an explicitly supplied dimension hypothesis. The record uses both ``edge-cone Yamabe metric'' and ``Yamabe minimizer''; the displayed statement follows the minimizer interpretation. Nonexistence of all constant-scalar-curvature edge-cone metrics would be a different, stronger assertion and is not established by this formulation.
\subsection{Short English statement}
For every \(1<\beta<2\), the conjecture asks whether the standard conformal class on \((S^n,S^{n-2})\), with cone angle \(2\pi\beta\), has no edge-cone metric attaining its Yamabe infimum. The dimension range, precise standard metric, and admissible behavior at the singular locus require clarification. The minimizer interpretation follows the supplied assessment.
\subsection{Sources}
\begin{itemize}
\item Local UnsolvedMath registry, \texttt{data/unsolvedmath/problems.json}, permanent ID \(30003545\), exact source problem code \texttt{catalogue entry}; supplied title: ``Nonexistence of Edge-Cone Yamabe Metrics at Small Cone Angles.'' No individual proposer is identified in the supplied record.
\item Supplied report citation: \emph{Analysis, Geometry and Topology of Positive Scalar Curvature Metrics} (2018), Oberwolfach Reports, pp.\ 2223--2298, DOI \texttt{10.4171/owr/2017/36}. Original OWR report: \url{https://doi.org/10.4171/owr/2017/36}. The record separately gives proposed year 2017.
\item Supplied literature reference: Akutagawa--Mondello, ``singular Yamabe minimizers,'' \url{https://arxiv.org/abs/1909.09367}. This descriptive reference label is reproduced from the record; no fuller bibliographic citation is supplied.
\item Literature assessment dated \texttt{2026-08-21}: the supplied curation labels the problem partially solved and asserts nonexistence of an edge-cone Yamabe minimizer for cone angle at least \(4\pi\), equivalently \(\beta\geq2\), while reporting the range \(2\pi<2\pi\beta<4\pi\) as unresolved. This is supplied dated curation, not an independently verified literature assessment.
\item The supplied statement-verification field reports that the official report was checked and that it proves the \(\beta\geq2\) result and states the remaining range as a conjecture. Its label \texttt{corrected\_verified} records supplied provenance; no source visit or independent verification is claimed here.
\end{itemize}
\end{document}
